\documentclass[11pt]{article}

\usepackage[a4paper,margin=1in]{geometry}
\usepackage{amsmath,amssymb,bm}
\usepackage{booktabs}
\usepackage{hyperref}

\hypersetup{
    colorlinks=true,
    linkcolor=blue,
    citecolor=blue,
    urlcolor=blue
}

\newcommand{\eps}{\varepsilon}
\newcommand{\mean}[1]{\langle #1 \rangle}
\newcommand{\master}{geqoe\_averaged\_zonal\_theory}
\newcommand{\secondorder}{geqoe\_averaged\_second\_order}
\newcommand{\tesseral}{geqoe\_tesseral\_thirdbody\_extension}

\title{Second-Order Tesseral Theory for GEqOE:\\
\large Resolving the Time-Angle Transcendental via
Fourier-Taylor Decomposition}
\author{Jose Manuel Montilla Garc\'ia}
\date{\today}

\begin{document}
\maketitle

\begin{abstract}
This note addresses a specific structural obstacle identified in the
tesseral extension plan (\texttt{\tesseral{}}): at second order in the
tesseral perturbation parameter, the frozen-$\theta_G$ approximation used
for first-order single-averaging breaks down, because the advance of the
Greenwich sidereal angle during one orbital revolution introduces a
transcendental factor $e^{-im\omega_E M(K)/n_0}$ that is not a rational
function of the complex true anomaly $F$. This breaks the
rational-in-$F$ structure required by the direct residue decomposition.

The resolution is a hybrid symbolic-numerical approach: the transcendental
time-angle factor is decomposed into a rapidly converging truncated
Fourier series (Jacobi-Anger expansion) in the eccentric anomaly $K$,
producing a finite Laurent polynomial in $z = e^{iK}$ and hence a
rational function of $F$. The convergence is controlled by the parameter
$x = m\omega_E g / n_0$, which is small ($x < 0.04$) for all practical
orbits, so that $N \lesssim 10$ Fourier terms suffice for machine
precision. The Fourier coefficients are Bessel functions $J_k(x)$,
computable once per orbit at negligible cost.

After this decomposition, the full second-order tesseral integrand
recovers the rational-in-$F$ structure and the direct residue
decomposition applies unchanged. The result is a hybrid theory:
symbolic in $(q, Q)$ (from the residue algorithm) and numerical in the
time-angle Fourier coefficients (from the Jacobi-Anger truncation).

This extension targets ultra-high-precision applications where the
$O(\eps_T^2) \sim 10^{-12}$ tesseral second-order terms become relevant,
such as precise orbit determination for geodetic satellites or
long-duration catalog maintenance.
\end{abstract}

\tableofcontents


%======================================================================
\section{The Problem}
\label{sec:problem}
%======================================================================

\subsection{Origin of the transcendental factor}

The first-order tesseral theory (\texttt{\tesseral{}}) freezes the
Greenwich sidereal angle $\theta_G$ during the one-revolution orbital
average. At first order, the freezing error is
$O(\eps_T \cdot \omega_E / n_0)$, which is second-order in the combined
perturbation counting and therefore negligible.

At second order in $\eps_T$, the freezing error becomes comparable to the
terms being captured. A rigorous $O(\eps_T^2)$ theory must account for
the advance of $\theta_G$ during the orbital average:
\begin{equation}
\theta_G(t(K)) = \theta_G(t_0) + \omega_E \cdot \frac{M(K)}{n_0},
\label{eq:theta_advance}
\end{equation}
where $t(K) - t_0 = M(K) / n_0$ is the time elapsed from the reference
epoch to the point $K$ on the orbit, and $M(K) = K - g\sin K$ is the
mean anomaly (generalized Kepler equation).

The tesseral forcing at the ``unfrozen'' level therefore contains factors
of the form
\begin{equation}
e^{-im\theta_G(t(K))}
= e^{-im\theta_{G0}} \cdot e^{-im\omega_E K / n_0}
\cdot e^{+im\omega_E g \sin K / n_0}.
\label{eq:unfrozen_phase}
\end{equation}

The first factor is a constant (absorbed into the epoch-dependent
coefficients). The second factor $e^{-im\omega_E K / n_0}$ is a complex
exponential at the non-integer frequency $m\omega_E / n_0$; it is
\emph{not} a Laurent polynomial in $z = e^{iK}$ when $m\omega_E/n_0
\notin \mathbb{Z}$ (which is the generic non-resonant case). The third
factor involves $e^{ix\sin K}$ with $x = m\omega_E g / n_0$, which is
also transcendental in $K$.

Together, these factors break the rational-in-$F$ structure that the
direct residue decomposition requires.

\subsection{Why this does not affect the first-order theory}

At first order, the tesseral short-period correction is obtained from:
\begin{equation}
n_0 \frac{\partial \bm u_T}{\partial K}
= \bm f_T(\bar{\bm z}, K, \theta_{G0})
- \mean{\bm f_T}_K,
\end{equation}
with $\theta_G$ frozen at $\theta_{G0}$. The frozen forcing
$\bm f_T(\bar{\bm z}, K, \theta_{G0})$ contains only the satellite's
orbital-geometry dependence on $K$, which is rational in $F$. The
transcendental factor \eqref{eq:unfrozen_phase} does not appear.

At second order, the effective forcing $\bm F_2$ (cf.\ the second-order
note, \texttt{\secondorder{}}, Eq.~17) contains products of first-order
quantities evaluated at the ``true'' phase $\theta_G(t(K))$, which is
where the unfrozen factor enters.


%======================================================================
\section{The Jacobi-Anger Resolution}
\label{sec:jacobi_anger}
%======================================================================

\subsection{Factoring the transcendental}

Rewrite \eqref{eq:unfrozen_phase} as
\begin{equation}
e^{-im\omega_E M(K)/n_0}
= e^{-im\omega_E K / n_0}
\cdot e^{+ix\sin K},
\qquad
x \equiv \frac{m\omega_E g}{n_0}.
\label{eq:factored_phase}
\end{equation}

The first factor $e^{-i\alpha K}$ with $\alpha = m\omega_E/n_0$ is a
pure complex exponential at a fixed (non-integer) frequency. It can be
handled by absorbing it into the integration variable or by treating
$\alpha$ as a parameter in the Fourier analysis.

The second factor admits the exact \textbf{Jacobi-Anger expansion}:
\begin{equation}
e^{ix\sin K}
= \sum_{k=-\infty}^{\infty} J_k(x)\, e^{ikK},
\label{eq:jacobi_anger}
\end{equation}
where $J_k(x)$ is the Bessel function of the first kind of order $k$.

\subsection{Rapid convergence}

The Bessel functions satisfy the asymptotic bound
\begin{equation}
|J_k(x)| \le \frac{1}{|k|!}\left(\frac{x}{2}\right)^{|k|}
\qquad \text{for } |k| \ge 1,
\label{eq:bessel_bound}
\end{equation}
which shows that the series \eqref{eq:jacobi_anger} converges
\emph{exponentially} in $|k|$ for any fixed $x$.

The parameter $x = m\omega_E g / n_0$ is small for all practical orbits:
\begin{center}
\begin{tabular}{lcccl}
\toprule
Regime & $\omega_E / n_0$ & $g$ & $x$ ($m=2$) & $N$ for $10^{-15}$ \\
\midrule
LEO circular   & 0.066 & 0.001 & $1.3 \times 10^{-4}$ & 2 \\
LEO moderate-$e$& 0.066 & 0.05  & $6.6 \times 10^{-3}$ & 3 \\
SSO            & 0.065 & 0.001 & $1.3 \times 10^{-4}$ & 2 \\
MEO/GPS        & 0.49  & 0.01  & $9.8 \times 10^{-3}$ & 4 \\
HEO/Molniya   & 0.046 & 0.74  & $6.8 \times 10^{-2}$ & 5 \\
GTO            & 0.052 & 0.73  & $7.6 \times 10^{-2}$ & 5 \\
GEO            & 1.003 & 0.001 & $2.0 \times 10^{-3}$ & 3 \\
\bottomrule
\end{tabular}
\end{center}

The column ``$N$ for $10^{-15}$'' gives the truncation order at which
$|J_k(x)| < 10^{-15}$ for all $|k| > N$, using the bound
\eqref{eq:bessel_bound}. In all cases, $N \le 5$ suffices for
double-precision machine epsilon.

\subsection{Truncated Laurent polynomial form}

Truncating \eqref{eq:jacobi_anger} at order $N$ and combining with the
first factor of \eqref{eq:factored_phase}:
\begin{equation}
e^{-im\omega_E M(K)/n_0}
\approx \sum_{k=-N}^{N} J_k(x)\,
e^{i(k - m\omega_E/n_0)K}.
\label{eq:truncated_phase}
\end{equation}

For the purpose of the residue decomposition, the non-integer frequency
shift $\alpha = m\omega_E / n_0$ can be handled in one of two ways:

\paragraph{Option A: Absorb into the integration contour.}
Change the averaging variable from $M$ to a shifted mean longitude
$\tilde M = M - \alpha M = (1 - \alpha)M$. This restores integer-frequency
exponentials but modifies the Jacobian.

\paragraph{Option B: Discrete Fourier projection.}
Instead of requiring exact periodicity in $K$, compute the Fourier
coefficients of the full integrand (including the $e^{-i\alpha K}$ factor)
numerically via discrete Fourier transform at $2N_{\mathrm{FFT}} + 1$
equally spaced points in $K$. This directly produces the Laurent
coefficients needed for the residue decomposition, bypassing the
non-integer frequency issue entirely.

\textbf{Option B is recommended} because it is fully general (works for
any non-autonomous perturbation, not just tesserals) and the FFT cost is
negligible ($\sim$microseconds).

\subsection{Recovery of rational-in-\texorpdfstring{$F$}{F} structure}

After the Fourier truncation (either option), the time-angle factor is
represented as a finite Laurent polynomial in $z = e^{iK}$:
\begin{equation}
e^{-im\omega_E M(K)/n_0}
\approx \sum_{k=-N_{\mathrm{tot}}}^{N_{\mathrm{tot}}}
b_k\, z^k,
\qquad
z = e^{iK}.
\label{eq:laurent_time_angle}
\end{equation}

Through the M\"obius map $F = (z - q)/(1 - qz)$, each monomial $z^k$ maps
to a rational function of $F$ with poles at $F = -q$ and $F = -1/q$
(from the denominator $(1-qz)^k$) and possibly at $F = 0$ (from
$z^k$ for $k < 0$). The pole \emph{locations} are the same as in the
zonal and first-order tesseral theories.

Therefore, the full second-order tesseral integrand---the product of
\eqref{eq:laurent_time_angle} with the rational-in-$F$ second-order
effective forcing---is a rational function of $F$ with poles only at
$F = 0$, $F = -q$, $F = -1/q$ (at higher multiplicities due to the
products). The direct residue decomposition algorithm applies
with the same modifications already identified in the second-order zonal
note (higher-order pole extraction).


%======================================================================
\section{Hybrid Symbolic-Numerical Architecture}
\label{sec:architecture}
%======================================================================

The resulting second-order tesseral theory has a natural two-layer
structure:

\subsection{Symbolic layer (offline, one-time)}

The residue decomposition is applied to the second-order tesseral
integrand with the time-angle Fourier coefficients $\{b_k\}$ carried as
\emph{symbolic parameters}. The output is a set of short-period kernel
expressions of the form
\begin{equation}
\bm u_{T,2}(F; q, Q, \{b_k\})
= \sum_k b_k \cdot \bm R_k(F; q, Q),
\end{equation}
where each $\bm R_k$ is a rational function of $F$ with coefficients in
$\mathbb{Q}(q, Q)$---exactly the same kind of expression produced by the
zonal residue decomposition.

The symbolic computation treats $\{b_k\}$ as free parameters and produces
$2N_{\mathrm{tot}} + 1$ kernel functions $\bm R_k$, one per Fourier mode.
This is a one-time cost, parallelizable over $k$.

\subsection{Numerical layer (once per orbit)}

For a specific orbit with parameters $(g, n_0)$:
\begin{enumerate}
\item Compute $x = m\omega_E g / n_0$.
\item Evaluate the Bessel coefficients $J_k(x)$ for $|k| \le N$.
\item Optionally, apply the FFT-based projection (Option~B) to obtain
the full set $\{b_k\}$ including the non-integer frequency shift.
\end{enumerate}

This takes microseconds and produces the numerical values of $\{b_k\}$
that are substituted into the symbolic expressions.

\subsection{Runtime evaluation}

At each output epoch, the second-order tesseral short-period correction is
\begin{equation}
\bm u_{T,2}(F)
= \sum_{k=-N_{\mathrm{tot}}}^{N_{\mathrm{tot}}}
b_k \cdot \bm R_k(F; q, Q),
\end{equation}
which is a weighted sum of $2N_{\mathrm{tot}} + 1$ rational function
evaluations. Since $N_{\mathrm{tot}} \lesssim 10$, this is comparable in
cost to evaluating $\sim$20 zonal short-period kernels---well within the
budget of any practical propagator.


%======================================================================
\section{Connection to Hansen Coefficients}
\label{sec:hansen}
%======================================================================

The Fourier decomposition of orbital quantities in the mean anomaly (or
eccentric anomaly) is closely related to the classical Hansen coefficients
$X_k^{n,m}(e)$, defined by
\begin{equation}
\left(\frac{r}{a}\right)^n e^{imf}
= \sum_{k=-\infty}^{\infty} X_k^{n,m}(e)\, e^{ikM}.
\end{equation}

The present approach differs from the classical Hansen machinery in three
respects:
\begin{enumerate}
\item The Fourier decomposition is applied to the \emph{time-angle
correction factor} $e^{-im\omega_E M/n_0}$ specifically, not to the
full orbital-geometry integrand. The orbital geometry is handled exactly
by the $F$-variable residue decomposition.

\item The Fourier coefficients (Bessel functions) are computed
\emph{numerically} for the specific orbit, not expanded as power series
in eccentricity. This avoids the slow convergence of Hansen coefficient
series for high eccentricity.

\item The hybrid approach keeps the $(q, Q)$ dependence in symbolic
(exact rational) form and isolates the time-angle coupling into a small
number of numerical coefficients. Classical Hansen-based theories carry
the eccentricity dependence in the Hansen coefficients themselves,
losing the exact symbolic structure.
\end{enumerate}

An alternative fully-symbolic route would use the Taylor expansion of the
Bessel functions:
\begin{equation}
J_k(x) = \sum_{j=0}^{\infty}
\frac{(-1)^j}{j!\,(j+k)!}\left(\frac{x}{2}\right)^{2j+k},
\end{equation}
with $x = m\omega_E g / n_0$, where $g$ is the GEqOE eccentricity
parameter. Truncating at a low order in $g$ would produce polynomial
expressions in $g$ for each $b_k$, restoring a fully symbolic theory at
the cost of an eccentricity expansion. For low-to-moderate eccentricity,
a handful of terms suffice; for high eccentricity, the numerical Bessel
evaluation is preferable.


%======================================================================
\section{Generalization via Taylor Algebra}
\label{sec:taylor_algebra}
%======================================================================

The Jacobi-Anger/Bessel approach is specific to the exponential of a
sinusoidal argument. A more general framework for handling arbitrary
smooth transcendental factors in the orbital average is provided by
\textbf{Taylor algebra} (differential algebra, DA), as implemented in
tools such as heyoka's jet transport or the COSY INFINITY framework.

\subsection{General procedure}

Given any smooth function $h(M(K))$ that is not a rational function of
$F$:
\begin{enumerate}
\item \textbf{Sample}: evaluate $h(M(K_j))$ at $2N_{\mathrm{FFT}} + 1$
equally spaced points $K_j = 2\pi j / (2N_{\mathrm{FFT}} + 1)$.
\item \textbf{FFT}: apply the discrete Fourier transform to obtain
Laurent coefficients $\{b_k\}_{|k| \le N_{\mathrm{FFT}}}$.
\item \textbf{Substitute}: use $\{b_k\}$ as numerical parameters in the
symbolic residue decomposition framework.
\end{enumerate}

This procedure works for \emph{any} smooth function of time along the
orbit, not just the tesseral time-angle factor. Potential applications
include:
\begin{itemize}
\item Atmospheric density models with complex altitude profiles
(transcendental in $K$ through $r(K)$).
\item Eclipse/shadow functions (non-smooth, but approximable by smooth
cutoffs).
\item Time-varying thrust profiles that depend on the orbital phase
through a non-polynomial control law.
\end{itemize}

\subsection{DA-enhanced Fourier coefficients}

For parametric studies where the Fourier coefficients must be known as
functions of orbital parameters (not just at a single orbit), Taylor
algebra provides automatic differentiation:
\begin{enumerate}
\item Initialize the orbital parameters $(g, Q, \omega_E/n_0, \ldots)$
as DA variables (Taylor jets).
\item Evaluate $h(M(K_j))$ at each sample point using DA arithmetic.
\item Apply the FFT to the DA-valued samples.
\item The resulting Fourier coefficients are themselves Taylor polynomials
in the orbital parameters, valid over a neighborhood.
\end{enumerate}

This produces a ``Taylor map of the Fourier coefficients''---a polynomial
approximation of the $b_k(g, \omega_E/n_0)$ dependence. Combined with
the symbolic-in-$(q, Q)$ kernel functions $\bm R_k$, this yields a
fully parametric second-order tesseral theory without ever solving the
transcendental issue in closed symbolic form.


%======================================================================
\section{Practical Relevance and When This Matters}
\label{sec:relevance}
%======================================================================

The second-order tesseral correction captures terms at
$O(\eps_T^2) \sim J_{2,2}^2 (R_e/a)^4 \sim 10^{-12}$ for LEO. This is:
\begin{itemize}
\item Below the first-order zonal truncation error
($O(\eps_Z^2) \sim 10^{-6}$) by six orders of magnitude.
\item Below the first-order tesseral correction itself
($O(\eps_T) \sim 10^{-6}$) by six orders of magnitude.
\item Below the zonal-tesseral cross term
($O(\eps_Z \eps_T) \sim 10^{-9}$) by three orders of magnitude.
\end{itemize}

Therefore, the second-order tesseral theory is only relevant when:
\begin{enumerate}
\item The second-order zonal theory is already implemented
($O(\eps_Z^2)$ captured).
\item The first-order tesseral and third-body theories are implemented
($O(\eps_T)$, $O(\eps_{3B})$ captured).
\item The cross-coupling layer is implemented
($O(\eps_Z \eps_T)$, $O(\eps_Z \eps_{3B})$ captured).
\item The application demands sub-millimeter or better analytical
accuracy (geodetic POD, gravity field recovery, inter-satellite ranging).
\end{enumerate}

The main value of this note is therefore \emph{architectural}: it
demonstrates that the rational-in-$F$ residue decomposition framework
is not fundamentally limited by the time-angle transcendental, and that
a clean resolution exists when the need arises.


%======================================================================
\section{Summary}
\label{sec:summary}
%======================================================================

\begin{enumerate}
\item The transcendental factor $e^{-im\omega_E M(K)/n_0}$ that obstructs
a purely rational-in-$F$ second-order tesseral theory is resolved by the
\textbf{Jacobi-Anger expansion}, which converts it to a rapidly converging
Fourier series in $K$.

\item The convergence parameter $x = m\omega_E g / n_0 < 0.08$ for all
practical orbits ensures that $N \le 5$ Fourier terms suffice for machine
precision.

\item After Fourier truncation, the second-order tesseral integrand
recovers the rational-in-$F$ structure and the \textbf{direct residue
decomposition applies unchanged}.

\item The resulting theory is \textbf{hybrid}: symbolic in $(q, Q)$
and numerical in the Bessel/Fourier coefficients $\{b_k\}$. The
numerical coefficients are computed once per orbit at negligible cost.

\item \textbf{Taylor algebra} (DA) generalizes the approach to arbitrary
smooth transcendental factors, providing automatic Fourier decomposition
and parametric sensitivity through jet transport.

\item The $O(\eps_T^2) \sim 10^{-12}$ magnitude of the second-order
tesseral terms makes this extension relevant only for ultra-high-precision
applications. Its primary value is demonstrating that the residue
decomposition framework has \textbf{no fundamental ceiling}---every layer
of the perturbation hierarchy can be captured within the same algorithmic
structure.
\end{enumerate}


\begin{thebibliography}{99}

\bibitem{bau2021}
Ba\`u, G., Hernando-Ayuso, J., Bombardelli, C. (2021).
A generalization of the equinoctial orbital elements.
\textit{Celestial Mechanics and Dynamical Astronomy}, 133, 50.

\bibitem{watson1944}
Watson, G.\,N. (1944).
\textit{A Treatise on the Theory of Bessel Functions}.
2nd ed., Cambridge University Press.

\bibitem{abramowitz1964}
Abramowitz, M., Stegun, I.\,A. (1964).
\textit{Handbook of Mathematical Functions}.
National Bureau of Standards.

\bibitem{berz1999}
Berz, M. (1999).
\textit{Modern Map Methods in Particle Beam Physics}.
Academic Press.

\bibitem{biscani2021}
Biscani, F., Izzo, D. (2021).
Revisiting high-order Taylor methods for astrodynamics and celestial
mechanics.
\textit{Monthly Notices of the Royal Astronomical Society}, 504(2),
2614--2628.

\bibitem{kaula1966}
Kaula, W.\,M. (1966).
\textit{Theory of Satellite Geodesy}.
Blaisdell, Waltham, MA.

\end{thebibliography}

\end{document}
